\documentclass[a4paper]{article}

\usepackage[fontset=fandol]{ctex}
\usepackage{amsmath, amssymb}
\usepackage{graphicx}
\usepackage[margin=2.5cm]{geometry}
\usepackage[most]{tcolorbox}
\usepackage{listings}
\usepackage{hyperref}
\usepackage{booktabs}
\usepackage{subcaption}
\usepackage{float}
\usepackage{tikz}

\setlength{\parindent}{2em}
\setlength{\parskip}{0.35em}
\sloppy

\newtcolorbox{knowledgebox}[1]{
    enhanced, colback=blue!5!white, colframe=blue!75!black, colbacktitle=blue!75!black,
    coltitle=white, fonttitle=\bfseries, title={#1}, attach boxed title to top left={yshift=-2mm, xshift=2mm},
    boxrule=1pt, sharp corners
}

\newtcolorbox{importantbox}[1]{
    enhanced, colback=yellow!10!white, colframe=yellow!80!black, colbacktitle=yellow!80!black,
    coltitle=black, fonttitle=\bfseries, title={#1}, sharp corners
}

\newtcolorbox{warningbox}[1]{
    enhanced, colback=red!5!white, colframe=red!75!black, colbacktitle=red!75!black,
    coltitle=white, fonttitle=\bfseries, title={#1}, sharp corners
}

\lstset{
    basicstyle=\ttfamily\small,
    keywordstyle=\color{blue},
    stringstyle=\color{red!60!black},
    commentstyle=\color{green!60!black},
    breaklines=true,
    frame=single,
    numbers=left,
    numberstyle=\tiny\color{gray},
    captionpos=b,
    extendedchars=false
}

\DeclareMathOperator*{\argmin}{arg\,min}
\newcommand{\HH}{\mathcal{H}}
\newcommand{\Dall}{\mathcal{D}_{all}}
\newcommand{\Dtrain}{\mathcal{D}_{train}}

\newcommand{\notetitle}{鱼与熊掌可以兼得的机器学习}
\newcommand{\notesubtitle}{李宏毅《机器学习》2025 · 为什么深度学习有潜力}
\newcommand{\noteauthors}{基于李宏毅老师课程整理}
\newcommand{\notedate}{\today}
\newcommand{\videochannel}{李宏毅深度学习课堂（Bilibili）}
\newcommand{\videoduration}{约 42 分钟}
\newcommand{\videourl}{https://www.bilibili.com/video/BV1TAtwzTE1S?p=11}

\begin{document}
\pagestyle{empty}

\begin{center}
    \vspace*{1.5cm}
    {\Huge\bfseries \notetitle\par}
    \vspace{0.4cm}
    {\LARGE \notesubtitle\par}
    \vspace{0.8cm}
    \includegraphics[width=0.62\textwidth]{video.jpg}\par
    \vspace{0.8cm}
    {\large \noteauthors\par}
    {\large \notedate\par}
    \vspace{0.8cm}
    \begin{tcolorbox}[width=0.82\textwidth,center,sharp corners,
                      colback=black!3!white,colframe=black!50!white]
        \centering
        \textbf{视频信息}\\[3pt]
        频道：\videochannel \\
        时长：\videoduration \\
        链接：\href{\videourl}{\videourl}
    \end{tcolorbox}
\end{center}

\newpage
\tableofcontents
\newpage
\pagestyle{plain}

% =====================================================================
\section{从 P08 的两难出发}
% =====================================================================

P08 的结尾留下一个两难。若模型家族 $\HH$ 很大，理想函数
\[
    h^{all}
    =
    \argmin_{h\in\HH}L(h,\Dall)
\]
有机会把理想 loss 压得很低；但 $\HH$ 大也会增加训练集与全体资料落差的风险。若 $\HH$ 很小，现实比较容易接近理想；但理想本身可能很差，因为候选函数太少。

这一讲要回答的问题是：有没有办法让理想 loss 低，同时又让模型复杂度不至于太大？这就是课堂说的「鱼与熊掌可以兼得」。

\begin{figure}[H]
    \centering
    \includegraphics[width=0.86\textwidth]{figures/fig01_tradeoff_goal.jpg}
    \caption{P08 的复杂度两难：大 $\HH$ 让理想好但现实可能远；小 $\HH$ 让现实近但理想可能差。目标是找到成员少、但成员足够好的 $\HH$。画面依据：约 00:00--02:25。}
    \label{fig:tradeoff_goal}
\end{figure}

课堂把理想状况说成：找到一个神奇的 $\HH$，它一方面候选函数不多，因此泛化风险较低；另一方面候选函数都是「精英」，能让 $L(h^{all},\Dall)$ 仍然很小。深度学习的潜力就在于，它可能用结构化的方式做到这件事。

\begin{importantbox}{这里的核心不是“模型越大越好”}
    如果只是把模型做得很大，当然可能让理想 loss 变低，但也会提高过拟合风险。深度学习真正有趣的地方是：对某些复杂但有规律的目标函数，深层结构能用比浅层结构少得多的参数表达它们。
\end{importantbox}

% =====================================================================
\section{复习：一个 Hidden Layer 已经能逼近任何函数}
% =====================================================================

在讲深度的好处前，老师先复习为什么 hidden layer 有用。考虑一维输入 $x$ 与输出 $y$，如果我们想逼近任意连续曲线，可以先用 piecewise linear function 近似它。切得越细，折线越接近原函数。

\begin{figure}[H]
    \centering
    \includegraphics[width=0.86\textwidth]{figures/fig02_piecewise_linear.jpg}
    \caption{任意曲线可以用足够多段 piecewise linear function 近似。画面依据：约 03:20--04:35。}
    \label{fig:piecewise}
\end{figure}

Piecewise linear function 又可以看成常数项加上一组阶梯状函数或折线函数。若 activation 是 sigmoid，可以用很陡的 sigmoid 近似 hard sigmoid；若 activation 是 ReLU，也可以用两个 ReLU 组合出类似 hard sigmoid 的形状。

\begin{figure}[H]
    \centering
    \includegraphics[width=0.86\textwidth]{figures/fig03_sum_of_steps.jpg}
    \caption{Piecewise linear function 可分解成常数项加上一组阶梯或折线形函数。画面依据：约 04:50--06:15。}
    \label{fig:sum_steps}
\end{figure}

\begin{figure}[H]
    \centering
    \includegraphics[width=0.86\textwidth]{figures/fig04_relu_hard_sigmoid.jpg}
    \caption{两个 ReLU 可组合出 hard sigmoid 形状；足够多 ReLU 可组合出足够多折线段。画面依据：约 07:30--08:40。}
    \label{fig:relu}
\end{figure}

因此，一个足够宽的一层 hidden network 可以逼近任意函数。这就是 universal approximation 的直觉版本。

\begin{knowledgebox}{那为什么还要 deep？}
    如果一个 hidden layer 已经能表示任何函数，问题就不再是“能不能表示”，而是“用多少参数表示”。深度学习的优势不是把不可表示的函数变成可表示，而是对很多有结构的函数，用更少参数、更有效率地表示。
\end{knowledgebox}

% =====================================================================
\section{实验观察：深一点通常比只变胖更有效}
% =====================================================================

课堂引用了一篇语音辨识论文的结果。随着层数增加，word error rate 下降：

\begin{center}
\begin{tabular}{cc}
\toprule
Layer $\times$ Size & Word Error Rate (\%)\\
\midrule
$1\times2k$ & 24.2\\
$2\times2k$ & 20.4\\
$3\times2k$ & 18.4\\
$4\times2k$ & 17.8\\
$5\times2k$ & 17.2\\
$7\times2k$ & 17.1\\
\bottomrule
\end{tabular}
\end{center}

\begin{figure}[H]
    \centering
    \includegraphics[width=0.86\textwidth]{figures/fig05_deeper_is_better.jpg}
    \caption{语音辨识实验中，层数加深后 word error rate 下降；但这还不能证明深度本身的价值，因为参数量也增加了。画面依据：约 13:30--14:55。}
    \label{fig:deeper_better}
\end{figure}

有人会质疑：深层网络表现好，只是因为参数更多、模型更大。若资料足够多，大模型当然能让理想更好；这不说明「深」特别重要。为了回答这个质疑，需要比较参数量相近的 shallow/fat network 与 deep/thin network。

\begin{figure}[H]
    \centering
    \includegraphics[width=0.86\textwidth]{figures/fig06_fat_vs_thin.jpg}
    \caption{Fat + Short vs. Thin + Tall：参数量相近时，把网络变深通常比单层变胖表现更好；一层 16k neuron 也不如较小的两层网络。画面依据：约 16:15--18:05。}
    \label{fig:fat_thin}
\end{figure}

实验给出的启发是：同样参数预算下，把参数沿深度方向组织起来，可能比全部塞进一层更有效。接下来要解释为什么。

\begin{warningbox}{不要把 deep learning 理解成“参数更多”}
    很多人把深度学习和大模型、大资料、容易 overfit 画等号。课堂要反过来强调：在某些目标函数上，deep network 可能用更少参数完成同样表达，因此反而更不容易 overfit，或需要更少训练资料。
\end{warningbox}

% =====================================================================
\section{为什么深层结构更有效？}
% =====================================================================

结论先说：若目标函数是 \textbf{复杂但有规律的}，deep network 往往优于 shallow network。复杂代表函数需要很多局部变化；有规律代表这些变化不是完全随机，而是可以通过组合结构重复利用。

\begin{figure}[H]
    \centering
    \includegraphics[width=0.86\textwidth]{figures/fig07_deep_more_effective.jpg}
    \caption{一层 hidden layer 可以表示任意函数，但深层结构通常更有效率。画面依据：约 18:20--20:10。}
    \label{fig:deep_effective}
\end{figure}

课堂用三个类比说明这种结构化效率。

\subsection{逻辑电路类比}

Parity check 的任务是判断输入 bit 序列里 1 的个数是奇数还是偶数。理论上，两层逻辑电路可以表示任何布尔函数，但如果强行用两层来做 parity check，可能需要 $2^d$ 级别的 gates，其中 $d$ 是输入长度。

更自然的做法是把 XNOR gates 串起来：前两个 bit 先合并，再和第三个合并，再和第四个合并。这样形成一个深的结构，使用的 gates 数量远少于两层暴力表示。

\begin{figure}[H]
    \centering
    \includegraphics[width=0.86\textwidth]{figures/fig08_logic_circuit_parity.jpg}
    \caption{逻辑电路类比：两层电路可表示任何布尔函数，但 parity check 用串接结构更有效率。画面依据：约 20:20--23:35。}
    \label{fig:logic}
\end{figure}

\subsection{程序设计类比}

写程序时，我们不会把所有逻辑塞进 main function，而会拆成 function、module、subroutine。这样同一功能可以反复调用，程序更短、更可读，也更容易维护。

神经网络中的深层结构也有类似意味：后层可以复用前层已经算出的中间表示，而不是每次从原始输入重新组合一切。

\subsection{剪窗花类比}

剪窗花时，直接在展开的纸上剪复杂图案很麻烦；先折纸，再剪几刀，展开后就得到复杂对称图案。折纸过程让简单操作通过结构重复产生复杂结果。

\begin{figure}[H]
    \centering
    \includegraphics[width=0.86\textwidth]{figures/fig09_papercut_analogy.jpg}
    \caption{剪窗花类比：折叠结构让少量剪刀动作产生复杂图案；深层网络也通过组合结构高效表达复杂函数。画面依据：约 24:45--25:45。}
    \label{fig:papercut}
\end{figure}

\begin{knowledgebox}{共同点}
    逻辑电路、程序模块、剪窗花的共同点是：有结构的组合能复用中间结果。Deep network 不是把参数随便堆高，而是让前层输出成为后层输入，让简单变换反复组合，产生指数级增长的表达片段。
\end{knowledgebox}

% =====================================================================
\section{一个 ReLU 构造：深层网络如何产生指数多片段}
% =====================================================================

课堂接着用一个一维 ReLU 网络展示深层结构的效率。先看一层网络，输入 $x$，两个 ReLU neuron：
\[
    a_1
    =
    \mathrm{ReLU}(x-0.5)
    +
    \mathrm{ReLU}(-x+0.5).
\]
它产生一个 V 字形关系。接第二层时，用同样结构作用在 $a_1$ 上，得到 $a_2$；由于 $a_1$ 本身已经是 V 字形，第二层会把它折成更多段，形成 W 字形。

\begin{figure}[H]
    \centering
    \includegraphics[width=0.86\textwidth]{figures/fig10_two_layers_more_pieces.jpg}
    \caption{两层 ReLU 结构把一层的 V 字形继续折叠，产生更多 piece。画面依据：约 29:50--33:05。}
    \label{fig:two_layers}
\end{figure}

继续叠第三层，会得到 $2^3$ 个 pieces。一般地，若有 $K$ 层，每层只放 2 个 ReLU neuron，总共 $2K$ 个 neuron，就能产生 $2^K$ 个线段片段。

若用 shallow network 做同样的锯齿状函数，一个 ReLU neuron 大致只能贡献一个新折点或一段变化；要产生 $2^K$ 个 pieces，就需要 $2^K$ 级别的 neuron。于是 deep 与 shallow 在参数需求上出现指数差距：
\[
    \text{deep: } O(K)
    \qquad
    \text{shallow: } O(2^K).
\]

\begin{figure}[H]
    \centering
    \includegraphics[width=0.86\textwidth]{figures/fig11_deep_vs_shallow_pieces.jpg}
    \caption{Deep network 用 $K$ 层、每层 2 个 neuron 产生 $2^K$ pieces；shallow network 若要产生同样 pieces，需指数多 neuron。画面依据：约 36:35--38:15。}
    \label{fig:pieces}
\end{figure}

\begin{importantbox}{为什么这能回应“鱼与熊掌”？}
    如果目标函数需要很多 pieces，但这些 pieces 有规律，deep network 可以用少量参数表达它。少量参数意味着候选函数集合更受限制，泛化风险较低；同时表达能力又足以让理想 loss 很低。因此深层结构有机会同时取得低理想 loss 与较小泛化落差。
\end{importantbox}

% =====================================================================
\section{复杂但有规律：哪些任务适合深度学习？}
% =====================================================================

深度的优势不是对所有函数都成立。如果目标函数完全随机，没有可复用结构，深层网络也未必占优。课堂强调的条件是：\textbf{复杂而有规律}。

语音、影像等任务很可能满足这个条件。它们的输入空间非常复杂，但其中有大量可组合结构：
\begin{itemize}
    \item 影像中，边缘组成纹理，纹理组成部件，部件组成物件。
    \item 语音中，短时声学特征组成音素，音素组成词，词组成句子。
    \item 语言中，字符或 token 组成词，词组成短语，短语组成语义结构。
\end{itemize}

这些层级结构正适合 deep network：低层提取局部简单结构，高层组合成更抽象的表示。

\begin{warningbox}{“深”不是魔法}
    深层结构只是在某些有组合规律的函数上更有效率。若任务不具备这种结构，或训练方式、资料量、优化方法不合适，深层网络仍然可能失败。深度学习强，是因为它的结构偏置常常贴合真实世界资料，而不是因为层数越多一定越好。
\end{warningbox}

% =====================================================================
\section{本讲综合：深度学习如何让鱼与熊掌兼得}
% =====================================================================

课堂最后给出结论：

\begin{figure}[H]
    \centering
    \includegraphics[width=0.86\textwidth]{figures/fig12_deep_learning_conclusion.jpg}
    \caption{结论：深度学习是一种让鱼与熊掌可以兼得的方法，有机会用较少参数得到较低 loss。画面依据：约 41:05--41:28。}
    \label{fig:conclusion}
\end{figure}

把这一讲和 P08 连起来，结论可以写成：
\[
    L(h^{train},\Dall)
    =
    L(h^{all},\Dall)
    +
    \bigl[
    L(h^{train},\Dall)-L(h^{all},\Dall)
    \bigr].
\]
第一项要靠表达能力压低；第二项要靠较小有效复杂度、足够资料和适当训练机制控制。浅层网络虽然理论上能表示任意函数，但表示某些复杂规律时可能需要指数多参数；参数一多，泛化就困难。深层网络通过组合结构，用较少参数表达同样复杂度的函数，因此有机会让两项同时变小。

\begin{knowledgebox}{一句话总结}
    深度学习的潜力不是“把模型做大”，而是“把模型做成深的组合结构”。当目标函数复杂但有规律时，深层结构能用更少参数表达它，于是既能让理想 loss 低，又能减轻过拟合压力。这就是这讲所说的鱼与熊掌可以兼得。
\end{knowledgebox}

\end{document}
